\chapter{Metodologia di confronto}
\label{cap:metodologia}

\begin{quote}\small
Questo capitolo descrive come i tre sensori sono stati confrontati. Fissa il
protocollo, le metriche e le convenzioni con cui sono stati acquisiti e letti i 441
trial del \cref{cap:confronto}, e presenta l'interfaccia web realizzata a supporto
della campagna, che risponde all'obiettivo O5. Gli strumenti software sono descritti
nel \cref{sec:piattaforma}. Qui si dice come sono stati usati.
\end{quote}

\section{Due strade complementari}
\label{sec:metodo-approccio}

Il confronto poggia su due attività distinte che si alimentano a vicenda. La prima
è sperimentale. I sensori e il soggetto vengono disposti in geometrie definite, e da
una prova all'altra cambia una sola variabile per volta. Può essere la distanza, la
quantità di movimento, il tipo di movimento, l'angolo o il materiale interposto.
Ogni configurazione è ripetuta e ogni ripetizione produce un file. La seconda
attività è numerica. I file formano un dataset omogeneo, nel formato unico del
\cref{sec:formato-dati}, analizzato a posteriori dalla catena del
\cref{sec:catena-analisi}. Le metriche non sono lette durante la prova ma calcolate
dopo, dagli stessi script e con le stesse convenzioni per tutti i sensori. Per questo
un risultato può essere ricalcolato e una figura rifatta senza ripetere
l'acquisizione, e ogni numero del \cref{cap:confronto} risale a un file e a una riga
di comando.

La divisione ha una conseguenza pratica. Durante la prova l'operatore non deve
leggere i numeri, gli basta sapere che i dati stanno arrivando e che i sensori
vedono quello che devono vedere. A questo serve l'interfaccia web della
\cref{sec:webui}, che mostra i sensori in tempo reale e permette di verificare il
posizionamento prima di registrare. Il suo ruolo nel confronto è di supporto e di
verifica, ed è dichiarato nella \cref{sec:webui-ruolo}.

\section{Protocollo sperimentale}
\label{sec:protocollo}
\label{sec:registro}

\paragraph{Setup.} Il \ldb\ e l'\pir\ sono letti nello stesso frame dal medesimo
firmware (\cref{sec:piattaforma}), sulla stessa base dei tempi. Il confronto non
richiede quindi alcun allineamento a posteriori. I due sensori stanno affiancati sul
bordo di un tavolo a circa 1~m di altezza, puntati verso l'area di prova, che è
interna al campo di entrambi. La stanza non contiene oggetti in movimento (nessun
ventilatore, tende ferme, porta chiusa). Il \ldc\ è stato provato in un secondo
tempo nelle stesse geometrie, con le deroghe dichiarate nella
\cref{sec:protocollo-ld2420}. Le quattro geometrie di prova sono schematizzate in
\cref{fig:geometrie}.

\begin{figure*}[htbp]
\forceversofloat
\centering
\includegraphics[width=\linewidth]{figures/fig27_geometrie}
\caption{Le geometrie della campagna, in scala. (a) La stanza, con i sensori sul
tavolo e le tacche a 1--5~m sull'asse. È la geometria di distanza, dose-risposta,
latenze, ostacoli e due persone. (b) Lo scenario del progetto, con i sensori
fissati sotto il piano e la persona rannicchiata a 50--60~cm. (c) L'arco a 1~m per
la copertura angolare. (d) Il corridoio della prova di portata, con i due vani porta
a 4 e 5~m attraverso cui sono visti i punti più lontani.}
\label{fig:geometrie}

\end{figure*}

\paragraph{Scenari.} Gli scenari sono raggruppati per domanda. La stanza vuota
misura i falsi positivi. La persona ferma a 2,3~m e sotto il banco a 60~cm misura i
falsi negativi nello scenario del progetto. La dose-risposta a 1~m cambia solo la
quantità di movimento, immobile, micro-movimenti e cammino sul posto. Il cammino sul
posto e l'attraversamento a 2--5~m cambiano solo il tipo di movimento. Le tacche a
1--5~m misurano la distanza. L'ingresso e l'uscita dalla stanza misurano le latenze.
L'arco a 1~m misura la copertura angolare, i pannelli interposti l'attenuazione, il
corridoio la portata massima, le due persone la capacità di separare i bersagli. Il
respiro a ritmo imposto usa gli stessi file e le stesse geometrie ed è trattato nel
\cref{cap:vitalita}. Ogni scenario è identificato da un nome, che è anche il nome dei
suoi file.

\paragraph{Ripetizioni e ground truth.} Ogni scenario è ripetuto cinque volte
(trial T01--T05), dieci per la latenza d'ingresso e tre per gli scenari angolari e a
due persone, che richiedono una seconda persona. I valori riportati nella forma
$\mu \pm \sigma$ sono media e deviazione standard su quelle ripetizioni. La presenza
reale e lo stato reale del soggetto sono dichiarati dall'operatore nei metadati di
ogni acquisizione (\cref{sec:formato-dati}). Per le latenze la ground truth statica
non basta, quindi il soggetto entra o esce a un istante annunciato a voce dal
calcolatore che registra. L'annuncio conta dalla prima riga di dati ricevuta, cioè
dalla stessa origine dei tempi usata dall'analisi. Un cronometro esterno
introdurrebbe un errore sistematico di 2--3~s, pari al riavvio del microcontrollore
all'apertura della porta seriale e maggiore della grandezza da misurare.

\paragraph{Configurazione fissa.} Il protocollo è stato definito prima delle
acquisizioni ed è rimasto invariato, perché una modifica a metà campagna avrebbe
reso i trial non confrontabili. Il PIR è stato caratterizzato nella configurazione a
lui più favorevole, con il ponticello di ritrigger in posizione H
(\cref{sec:interpretazione-pir}). Le soglie del \ldb\ sono quelle di fabbrica per
tutta la campagna (\cref{sec:soglie}). Il tipo di movimento è annotato insieme alla
distanza, perché il PIR risponde al primo e non alla seconda (\cref{sec:dose-risposta}).

\paragraph{Registro.} Ogni sessione è annotata in un registro con data, ora,
temperatura della stanza, soggetto, alimentazione e note. Le sessioni si collocano
fra 25 e 29\,\textdegree C, una condizione sfavorevole al PIR per il rilevamento
(\cref{sec:pir-temperatura}), di cui la discussione tiene conto. Il registro completo
è nel repository del lavoro (\cref{app:codice}).

\section{Metriche e convenzioni}
\label{sec:metriche}

\paragraph{Metriche.} Le metriche sono calcolate dallo script di analisi del
\cref{sec:catena-analisi} su ogni trial e poi aggregate per scenario.
\begin{description}[leftmargin=1.4em,style=nextline,itemsep=2pt]
  \item[Falsi negativi (\%)] campioni con presenza reale in cui il sensore riporta
    assenza. È la metrica dello scenario della persona ferma. Il suo complemento, la
    frazione di campioni con presenza riportata, è il tasso di rilevamento.
  \item[Falsi positivi (eventi/h)] transizioni da assenza a presenza per ora di
    stanza vuota. Si contano eventi e non campioni, perché un falso positivo di
    dieci secondi non è dieci volte più grave di uno di un secondo.
  \item[Latenza di rilevamento e di rilascio] tempo fra l'evento e la prima
    transizione stabile del sensore, all'ingresso e all'uscita. Per il confronto
    fra sensori si usa la differenza appaiata sullo stesso trial, in cui tragitto e
    tempo di reazione si cancellano.
  \item[Accuratezza della distanza] scostamento fra distanza riportata dal radar e
    distanza reale, separando la dispersione fra trial da quella entro il trial.
  \item[Energia] valore normalizzato 0--100 riportato dal \ldb\ per il bersaglio e
    per gate. Non è una potenza e non si converte in decibel. Serve come grandezza
    graduata del movimento e come testimone dello stimolo, perché due prove con la
    stessa energia hanno riprodotto lo stesso movimento.
\end{description}

\paragraph{Zero eventi.} Quando in una finestra di osservazione non si verifica
alcun evento non si scrive zero eventi all'ora. Si riporta il limite superiore dato
dalla regola del tre, $3/T$ con $T$ in ore, che approssima il limite al 95\,\% di
confidenza per un processo di Poisson senza eventi. Il numero dipende dal tempo di
osservazione e non dal sensore.

\paragraph{Transitorio iniziale.} I primi secondi di ogni trial sono scartati,
perché contengono il posizionamento del soggetto e la coda di presenza del radar.
La durata dello scarto dipende dallo scenario ed è la stessa in tutti gli script e
in tutte le figure. Vale 20~s negli scenari ordinari e 40~s sotto il banco, dove il
soggetto deve anche posizionarsi. Vale 60~s a stanza vuota, perché l'operatore deve
uscire, 120~s negli scenari di selettività, dove la coda del radar dopo il
posizionamento arriva a superare i 100~s (\cref{sec:angolare}), e 90~s nei trial da
fermo del \ldc, che rilascia in circa 55~s.

\paragraph{Validità dei file.} Un file entra nelle statistiche solo se è un trial.
I controlli del fondo, le prove tecniche e i file marcati non validi nel registro
sono elencati in una lista di esclusione degli script, con il motivo, e non entrano
in nessuna tabella. Un file non valido è tenuto nel dataset e rinominato, non
cancellato. Il caso tipico è il radar muto per un cavo mosso, riconoscibile da
centinaia di frame identici.

\section{Protocollo per il HLK-LD2420}
\label{sec:protocollo-ld2420}

L'esemplare di \ldc\ in dotazione vede una persona solo fino a circa 2~m contro gli
8 dichiarati, per un difetto attribuito con ogni probabilità al trasmettitore
(\cref{sec:ld2420-stato}). La campagna sul secondo radar è stata quindi svolta per
intero entro i 2~m, ripetendo i test del \ldb\ che non richiedono
distanze maggiori, per 161 trial validi. Restano fuori solo le prove che le
richiedono. Entrano invariate le due più importanti per il progetto, la dose-risposta
a 1~m e la persona sotto il banco.

Quattro deroghe alla parità di configurazione vanno dichiarate. I dati sono letti
nella modalità binaria ricostruita dalla comunità (\cref{sec:ld2420}), l'unica che
esponga le energie per gate. Le soglie sono tarate sul fondo misurato e il gate
massimo è 6, perché a configurazione di fabbrica il modulo non rilascia mai la
presenza nella stanza di prova (\cref{sec:ld2420-risultati}). Dal 7 settembre il
gate~0 è portato sopra la coda del rumore, in quella che il testo chiama
configurazione v2. Il PIR non è cablato nei test del solo \ldc, e il firmware scrive
il valore sentinella al posto del suo stato (\cref{sec:formato-dati}), così che
nessuna metrica del PIR venga calcolata su quei file. Ogni numero del \ldc\ è
attribuito all'esemplare, mai al modello.

Due accorgimenti operativi completano il protocollo. Dopo ogni accensione nessuno
resta entro 2~m dal modulo per 90~s, perché il fondo viene stimato all'avvio e un
corpo vicino lo corrompe. Prima di ogni sessione un minuto di stanza vuota verifica
che il fondo per gate sia quello atteso.

\section{L'interfaccia web di supporto}
\label{sec:webui}
\label{cap:webui}

L'obiettivo O5 chiede che l'ESP32 pubblichi i dati, che una pagina web li mostri in
tempo reale, li salvi con le statistiche di sessione e li esporti in CSV. Il
risultato è una dashboard servita direttamente dal nodo, senza alcuna rete esterna,
che produce lo stesso formato dati della catena seriale. Questa sezione ne riassume
requisiti, progetto, realizzazione e verifica, e dichiara il ruolo che ha avuto nella
campagna.

\subsection{Requisiti e architettura}
\label{sec:webui-requisiti}

Dall'obiettivo discendono cinque requisiti funzionali e due vincoli. L'ESP32 deve
rendere disponibili i dati in rete (R1). La visualizzazione deve seguire il sensore
alla cadenza dell'acquisizione, 5~Hz (R2). La stessa applicazione deve calcolare le
statistiche di sessione (R3) ed esportare i dati in CSV (R4). Deve mostrare l'indice
di vitalità calcolato a bordo (R5, \cref{cap:vitalita}). I due vincoli vengono dal
contesto. In emergenza sismica non c'è connettività, né verso Internet né verso una
rete locale, quindi tutto deve funzionare senza alcuna rete esistente (N1). Il CSV
esportato dal browser deve avere lo stesso formato di quello prodotto dalla catena
seriale, così che in tutta la tesi esista un solo formato dati (N2). Le tre
architetture considerate sono in \cref{tab:webui-opzioni}.

\begin{table*}[htbp]
\forcerectofloat
\centering
\caption{Opzioni architetturali valutate per la web UI. È stata scelta la A.}
\label{tab:webui-opzioni}

\small
\begin{tabularx}{\linewidth}{@{}lXX@{}}
\toprule
\textbf{Architettura} & \textbf{Vantaggi} & \textbf{Svantaggi} \\
\midrule
\textbf{A. Sito servito dall'ESP32} (web server e WebSocket a bordo) &
  nessuna infrastruttura, funziona offline, coerente con lo scenario di emergenza &
  RAM e flash limitate, pochi client simultanei \\
B. ESP32, broker MQTT, server con database &
  scalabile, storico illimitato, più vicino a una piattaforma reale &
  richiede broker e server sempre accesi, infrastruttura eccessiva per una tesi di
  sensing \\
C. Piattaforma cloud (ThingSpeak, Blynk) &
  pronta all'uso &
  non funziona offline, cadenza limitata a circa 15~s, dati fuori controllo \\
\bottomrule
\end{tabularx}
\end{table*}

È stata scelta l'opzione~A, confermata dal relatore. Un nodo che serve la propria
dashboard senza alcuna rete esterna è un caso concreto di scenario senza
connessione, che è lo scenario del progetto. L'opzione~C viola il vincolo N1, e la~B
lo rispetta solo con un server locale che in emergenza non ci sarebbe. Il volume dei
dati è piccolo, una trentina di valori a 5~Hz, circa 2~kB/s. La persistenza non
serve a bordo, perché lo storico di sessione si accumula nel browser ed è il browser
a produrre il CSV. L'ESP32 crea quindi la propria rete WiFi in modalità Access Point
e non tenta mai di collegarsi ad altre. L'opzione~B resta documentata come
alternativa valutata e scartata.

\subsection{Progetto e realizzazione}
\label{sec:webui-progetto}
\label{sec:webui-piano}
\label{sec:webui-lezioni}
\label{sec:webui-riferimenti}

\paragraph{Firmware.} La web UI è il logger del \cref{sec:piattaforma} con in più la
rete. Stessa libreria per il radar (\texttt{MyLD2410}~\cite{myld2410}), stesso
cablaggio, stessa cadenza di 200~ms e stesso CSV sulla seriale, nel formato a strati
del \cref{sec:formato-dati}. Il server HTTP e WebSocket è la libreria ESP Async
WebServer nel fork mantenuto ESP32Async~\cite{esp32async}, necessario perché
l'originale non compila con il core ESP32 corrente. La versione asincrona serve le
richieste da un task sull'altro core, mentre il ciclo principale continua a leggere
il radar a 5~Hz. La pagina è incorporata nel firmware come costanti compresse in
flash, circa 80~kB, e richiede lo schema di partizione a 3~MB per l'applicazione,
l'unica impostazione da cambiare rispetto al logger. Un piccolo server DNS a bordo
risponde con l'indirizzo del nodo a qualunque nome richiesto, così la dashboard si
apre da sola collegandosi alla rete del nodo, senza conoscere un indirizzo IP.

\paragraph{Pagina.} La pagina è in JavaScript puro, tre file per circa 800 righe,
senza framework né passo di build. Il tempo reale lo fa il canale WebSocket,
disponibile in ogni browser senza librerie. L'unica dipendenza è la libreria di
grafici Chart.js~\cite{chartjs}, servita dall'ESP32 insieme alla pagina. Come
riferimento di interfaccia è stato esaminato un configuratore web open source per
il \ldb~\cite{ld2410configurator}, da cui è ripresa la rappresentazione a barre delle
energie per gate con la soglia sovrapposta. La pagina ha quattro aree. Lo stato
istantaneo dei due sensori. La gauge dell'indice di vitalità, con la convenzione a
semaforo in ordine inverso all'indice, rosso per la classe bassa che per il soccorso
è la più urgente (\cref{sec:vitalita-classi}). I tre grafici degli ultimi 60~s, con
l'energia, le barre per gate con le soglie lette dal modulo all'avvio e la timeline
della presenza di radar e PIR sulla stessa base dei tempi. La sessione di
registrazione, con gli stessi metadati sperimentali della catena seriale. Le soglie
non sono modificabili dalla pagina, deliberatamente, perché devono restare fisse per
tutta la campagna (\cref{sec:soglie}).

\paragraph{Protocollo dati.} A ogni campione il nodo invia ai client il messaggio del
\cref{lst:ws}, a 5~Hz. Contiene tutte le grandezze del CSV del logger più le due
calcolate a bordo, così che il CSV esportato dal browser sia equivalente a quello
della seriale (N2). Il nodo conserva uno storico di 60~s e lo invia a ogni client che
si collega, così il grafico è pieno anche dopo un aggiornamento della pagina. Le
statistiche sono calcolate nel browser una volta al secondo. Il CSV esportato ha le
colonne del logger e i metadati dello script di acquisizione, nello stesso ordine,
più due colonne in coda con l'indice e la classe calcolati dal firmware. Stanno in
coda perché gli script esistenti, che leggono le colonne per nome, le ignorino senza
modifiche. Sono il dato usato dalla verifica del porting (\cref{sec:vitalita-bordo}).

\begin{lstlisting}[caption={Messaggio WebSocket inviato a 5~Hz a ogni client collegato.},
label={lst:ws},basicstyle=\ttfamily\scriptsize]
{
  "t": 123456,                        // millis() dell'ESP32
  "radar_ok": 1,                      // frame valido ricevuto in questo periodo
  "presence": 1, "moving": 1, "still": 0,
  "mdist": 145, "sdist": 0,           // distanze in cm
  "menergy": 72, "senergy": 0,        // energia del bersaglio 0-100
  "pir": 1,                           // sensore PIR, stesso frame
  "gates_m": [0,5,60,10,0,0,0,0,0],   // energia per gate, movimento
  "gates_s": [0,0,45,8,0,0,0,0,0],    // energia per gate, stazionario
  "light": 24, "out": 1,              // fotosensore e stato del pin OUT
  "vitality": 54,                     // indice di vitalita' 0-100
  "vitality_class": "vitalita_moderata"
}
\end{lstlisting}

\paragraph{Realizzazione.} Lo sviluppo è stato organizzato in cinque passi, ciascuno
con un criterio di accettazione verificabile prima di passare al successivo. Rete e
pagina statica, con la pagina che si apre da telefono e da portatile. Lettura dei
sensori e WebSocket, con i valori in pagina che seguono il sensore senza accumulo di
latenza. Grafici e storico, con le barre per gate che indicano il gate corretto
spostando il soggetto di 0,75~m. Sessione ed esportazione, con il CSV del browser
identico a quello della seriale. Indice di vitalità a bordo, uguale a quello
calcolato in Python entro $\pm 1$ a regime (\cref{sec:vitalita-bordo}). Un problema
incontrato vale per chi ripete il lavoro. I callback del server asincrono girano su
un task diverso dal ciclo principale, e una stampa di log alla connessione di un
client ha corrotto una riga CSV in scrittura. Gli eventi di rete passano ora da una
coda e vengono stampati dal ciclo principale fra una riga e l'altra.

\subsection{Verifica di equivalenza: un solo formato dati}
\label{sec:webui-verifica}

Il vincolo N2 lega l'interfaccia al resto della tesi, e la sua verifica è stata
svolta come test di equivalenza. La stessa sessione è stata registrata in due modi
insieme, dal browser con la funzione di esportazione e dalla seriale con lo script di
acquisizione lanciato in parallelo sulla porta del nodo. I due file sono stati
confrontati campione per campione allineandoli sul tempo dell'ESP32. Sui 1346
campioni comuni, cioè l'intera sovrapposizione temporale, tutte le colonne sono
identiche in ogni valore, energie per gate e distanze comprese. Lo script di analisi
legge entrambi i file e produce le stesse metriche. Non sono formati compatibili, è
lo stesso dato raggiungibile per due vie.

\subsection{Ruolo nella campagna}
\label{sec:webui-ruolo}

L'interfaccia è arrivata nella seconda metà della campagna, e il suo ruolo va
dichiarato per quello che è stato. Il dataset comparativo del \cref{cap:confronto} è
acquisito dalla seriale con lo script di acquisizione, che esisteva prima
dell'interfaccia. La web UI ha avuto tre funzioni. È stata il monitor in tempo reale
durante il posizionamento dei sensori e del soggetto, in particolare in corridoio,
dove la timeline della presenza di radar e PIR mostra dal vivo il risultato del
confronto. È stata il canale di acquisizione della prova di portata
(\cref{sec:portata-massima}), resa possibile dalla verifica di equivalenza senza
alcuna modifica alla catena di analisi. È stata infine la sede del calcolo a bordo
dell'indice di vitalità e del suo confronto con il riferimento offline
(\cref{sec:vitalita-bordo}). Non è stata usata per calcolare le metriche di
confronto, che restano compito della catena di analisi sui file.
